
This paper measured benchmark score compression at field scale, finding that 31.1\% of 466 qualifying ML leaderboard benchmarks exhibit at least one compression event over a 7-year panel (2018–2025). Granger causality analysis provided evidence that submission count accumulation temporally precedes score variance compression at lag 2 (p = 1.854 × 10⁻⁵), with reverse causality not confirmed. Domain-specific health estimators showed large discriminative effect sizes on synthetic panels (|d| > 5 in all domains), and NLP H_d achieved AUC = 0.857 for cross-sectional identification of compressed benchmarks on real data.

The prospective prediction component of the BCBHS framework — Cox proportional hazards survival modeling (H-M3) and Kaplan-Meier lead-time analysis (H-M4) — was not executed. The pre-registered collapse event operationalization (expanding Kendall τ > 0.85 on quarterly score variance) yielded only 1 event due to a structural incompatibility with the compression signal defined in H-M1, blocking the survival model pipeline. A persistence-based collapse criterion (≥ 4 consecutive compressed quarters) is proposed as the resolution.

The principal contributions of this work are: (1) the first systematic cross-domain measurement of benchmark compression prevalence; (2) Granger-causal evidence for the temporal structure of the submission-to-compression mechanism; (3) validated cross-sectional discriminability for the NLP H_d signal on real data; and (4) principled falsification of the expanding-τ collapse operationalization with a concrete resolution path. The framework's central prediction claim remains to be validated.

